\section{The simulation argument in philosophy and formal epistemology}
\label{sec:philosophy}

Most scientific interest in whether the universe is a computation goes back to a
probabilistic argument by Bostrom~\cite{Bostrom2003are}. This section states that
argument exactly, sets out the main lines of objection in the philosophical literature,
and says what the argument does and does not imply for empirical work. That last point
connects this section to the physical proposals reviewed later.

\subsection{Bostrom's argument}
\label{sec:philosophy:bostrom}

\paragraph{Premises.}
The argument has one metaphysical premise, one empirical premise and one epistemic
premise. The metaphysical premise is a weak \emph{substrate-independence} thesis. Bostrom
does not require that it be necessarily true, only that ``in fact, a computer running a
suitable program would be conscious''. He needs only that replicating a brain's
computational processes ``in suitably fine-grained detail, such as on the level of
individual synapses'' would suffice for experience. He takes this ``as a given'' rather
than arguing for it~\cite[\S\,II]{Bostrom2003are}.

The empirical premise concerns computing capacity. Bostrom estimates the processing needed
to emulate one human brain at ${\sim}10^{14}$ or ${\sim}10^{16}$--$10^{17}$ operations
per second. He estimates the processing needed for an \emph{ancestor-simulation} of the
mental history of humankind at roughly $10^{33}$--$10^{36}$ operations. The second figure
is the product of $10^{11}$ humans, 50 years each, $3\times10^{7}$\,s per year and
$10^{14}$--$10^{17}$ operations per second~\cite[\S\,III and fn.~10]{Bostrom2003are}. He
takes $10^{42}$ operations per second as a rough figure for a planetary-mass computer
built from known nanotechnological designs. On that figure, one such computer could run
an ancestor-simulation ``by using less than one millionth of its processing power for
one second''~\cite[\S\,III]{Bostrom2003are}. The $10^{42}$ figure comes from an unpublished
working manuscript that Bostrom cites. It is not derived from physical law.

The resource estimate depends on a further assumption: a simulation need only reproduce
what the simulated observers can notice. Detail is filled in ``on an as-needed basis'',
and anomalies are handled by editing the brain states of the observers who noticed them
or by re-running the simulation~\cite[\S\,III]{Bostrom2003are}. Bostrom himself states that
simulating ``the entire universe down to the quantum level is obviously infeasible,
unless radically new physics is discovered''~\cite[\S\,III]{Bostrom2003are}.

\paragraph{The fraction formula.}
Let $f_P$ be the fraction of human-level technological civilizations that reach a
posthuman stage. Let $\bar N$ be the average number of ancestor-simulations run by a
posthuman civilization, and $\bar H$ the average number of individuals who lived in a
civilization before it became posthuman. Bostrom writes the fraction of observers with
human-type experiences who live in simulations as~\cite[\S\,IV]{Bostrom2003are}
\begin{equation}
  f_{\mathrm{sim}} \;=\; \frac{f_P\,\bar N\,\bar H}{(f_P\,\bar N\,\bar H) + \bar H}.
  \label{eq:philosophy:fsim}
\end{equation}
Let $f_I$ be the fraction of posthuman civilizations interested in running
ancestor-simulations, and $\bar N_I$ the average number such civilizations run. Then
$\bar N = f_I \bar N_I$ and
\begin{equation}
  f_{\mathrm{sim}} \;=\; \frac{f_P\, f_I\, \bar N_I}{(f_P\, f_I\, \bar N_I) + 1}.
  \label{eq:philosophy:fsimstar}
\end{equation}
Bostrom argues that $\bar N_I$ is extremely large, and concludes that at least one of
the following holds~\cite[\S\,IV]{Bostrom2003are}:
(1)~$f_P \approx 0$; (2)~$f_I \approx 0$; (3)~$f_{\mathrm{sim}} \approx 1$.

\paragraph{The indifference step.}
A separate premise turns proposition~(3) into a credence. This is the \emph{bland
indifference principle},
$\mathrm{Cr}(\mathrm{SIM} \mid f_{\mathrm{sim}} = x) = x$. Bostrom says it ``prescribes
indifference only between hypotheses about which observer you are, when you have no
information about which of these observers you are''. He supports it with an analogy to
an unexpressed junk-DNA sequence and a betting argument~\cite[\S\,V]{Bostrom2003are}. He
also says it is weaker than the premise of the Doomsday argument, which he describes as
``much stronger and more controversial''~\cite[\S\,V]{Bostrom2003are}.

\paragraph{What Bostrom claims.}
The conclusion is a disjunction, and Bostrom has repeatedly declined to endorse its third
disjunct on its own. In 2003 he recommended apportioning credence ``roughly evenly''
among (1)--(3)~\cite[\S\,VII]{Bostrom2003are}. In 2005 he wrote that ``the disjunction is
all that the simulation argument purports to show''~\cite[\S\,1]{Bostrom2005simulation}. In
2009 he wrote: ``I believe that we are probably not simulated''~\cite{Bostrom2009simulation}.
His current FAQ assigns a ``substantial probability'' to the hypothesis but gives no
number~\cite[Q2]{Bostrom2025simulation}. Bostrom also held that the truth of~(3) would
leave most beliefs nearly intact, since ``our best guide'' to how the simulators set up
our world ``is the standard empirical study of the universe we
see''~\cite[\S\,VI]{Bostrom2003are}.

\paragraph{The 2011 patch.}
Bostrom and Kulczycki later identified ``a mathematical non sequitur'' in
Eq.~\eqref{eq:philosophy:fsim}~\cite{Bostrom2011patch}. The formula uses one average
$\bar H$ for all civilizations. Suppose instead that civilizations which go on to run
simulations have much shorter pre-posthuman phases than those that do not. Then all
three disjuncts can fail at once. The paper offers two independent patches. The first
assumes that cumulative pre-posthuman populations do not differ by an astronomically
large factor between simulating and non-simulating civilizations; its appendix uses a
factor of one million. The second restricts the reference class using indexical evidence
about our position in history, such as a ``computer age birth rank''. Throughout, the
paper assumes a finite universe~\cite{Bostrom2011patch}.

\subsection{Objections to the inference}
\label{sec:philosophy:objections}

\paragraph{The indifference step and the nature of evidence.}
Weatherson~\cite{Weatherson2003are} distinguishes four readings of the general principle
behind the indifference step. On two readings, he argues, it is false. On the third it
does not yield Bostrom's principle. On the fourth it yields the principle only given an
auxiliary premise: that all of one's evidence is probabilistically independent of being
a Sim, which ``we have no reason to believe''. Against that premise he raises
epistemological externalism: ``seeing a dagger and sim-seeing a sim-dagger'' may give
different evidence even if the experiences match. He also raises a Goodman-style
challenge built on the gerrymandered predicate ``suman''~\cite{Weatherson2003are}.

Bostrom replied that the principle constrains \emph{conditional} credence, so that
further evidence enters by conditionalization. On his reply, the argument requires only
that our actual evidence is too weak to move credence far once
$f_{\mathrm{sim}}\approx 1$ is given~\cite[\S\,2]{Bostrom2005simulation}. He also argued
that the simulation argument differs from brain-in-a-vat scepticism. It relies on
empirical premises and aims to raise credence in a disjunction about the world, not to
produce general agnosticism~\cite[\S\,4]{Bostrom2005simulation}.

Birch~\cite{Birch2013on} makes explicit the constraint that the argument places on
credence:
$\mathrm{Cr}(\mathrm{SIM}) \ge \mathrm{Cr}(\mathrm{SIM}\mid H_3)\,
\bigl(1-\mathrm{Cr}(H_1\vee H_2)\bigr)$.
He then argues that the argument needs a ``selective scepticism''. The trilemma requires
good evidence about the physical limits of computation. The indifference step requires
that our evidence not support mundane claims about our own constitution. Given a
``parity of evidence'' premise, these requirements are ``jointly incompatible''.
Retreating to a weaker four-way disjunction does not recover the constraint on credence.
Birch also argues that a simulated observer could infer with confidence only that at
least one ancestor-simulation is possible. He concludes that ``there is thus no good
reason to believe'' the argument's conclusion~\cite{Birch2013on}.

Crawford (abstract-level) gives a structurally parallel argument that one is a ``freak
observer'' formed in a fluctuation. He also argues that ``the evidence that simulants or
freak observers exist is not a reason to think that I am one of them''~\cite{Crawford2013freak}.

Thomas~\cite{Thomas2026simulation} accepts that Bostrom's reasoning establishes only a
conditional claim. He notes that confidence in a high ratio of simulants is hard to
square with being a simulant. He therefore replaces that premise with a weaker one: the
expected simulant-to-non-simulant ratio in one's reference class is high,
\emph{conditional on} one's being non-simulated. Given an indifference principle and an
admissibility condition on the reference class, he proves that the odds of being
simulated are at least this expected ratio. He also argues that finding many simulations
being run would not necessarily raise those odds~\cite{Thomas2026simulation}. This
runs against Bostrom's remark that our going on to create ancestor-simulations ``would be
strong evidence against (1) and (2)''~\cite[\S\,VI]{Bostrom2003are}.

\paragraph{Substrate-independence.}
Some objections target the metaphysical premise. Beisbart (abstract-level) argues that
the interactions among hardware components during a brain simulation ``do not resemble
the ways in which neurons interact in the brain''~\cite{Beisbart2014are}. Summers and
Arvan (abstract-level) argue that if panpsychism or panqualityism is true, living in a
simulation might be possible only as a brain in a vat, which would make simulation
unlikely~\cite{Summers2022two}. Chalmers reports that Godfrey-Smith~\cite{GodfreySmith2024simulation} favours a
biological basis for experience, and replies that functionalist views remain serious
possibilities~\cite{Chalmers2024taking}.

\paragraph{The motivational disjunct.}
Bostrom noted that an ethical ban on ancestor-simulations would support~(2) only if it
were combined with a civilization-wide means of enforcing it~\cite[\S\,VI]{Bostrom2003are}.
Greene~\cite{Greene2020termination} argues that running ancestor-simulations, and probing
experimentally for simulation, carry a risk that the simulators terminate our world.
Harris~\cite{Harris2024simulation} builds on this. He argues that any beings unsure
whether they are simulated have self-preservation reasons not to run conscious
simulations, since a nested simulation could exhaust the resources of a lower level.
This supports disjunct~(2), and Harris concludes that a high probability for~(3) ``is
premature''~\cite{Harris2024simulation}.

\subsection{Nesting and resource decay}
\label{sec:philosophy:nesting}

Bostrom allowed that simulated civilizations might run simulations of their own. He also
noted that the cost to the basement level counts against many levels, so that ``we
should expect our simulation to be terminated when we are about to become
posthuman''~\cite[\S\,VI]{Bostrom2003are}. Brueckner objected that a simulated being cannot
``really'' build another computer~\cite{Brueckner2008simulation}. Our knowledge of this
objection comes from the opening paragraph of Brueckner's paper and from Bostrom's
quotation of it. Bostrom replied that virtual machines really perform the computations they
emulate~\cite{Bostrom2009simulation}.

Carroll~\cite{Carroll2016maybe} states a ``resolution conundrum''. Each level has less
computing power than the one above it, so the hierarchy ends in simulations unable to
simulate. With a typicality premise, most observers then live at that lowest level. This
conflicts with the premise that such simulations are feasible. Carroll calls the
typicality premise ``more or less completely without justification''. Kipping's
hierarchical model agrees that most realities lie in the lowest level when $p\lambda\gg1$
(notation below). He suggests that inhabitants of the lowest level might still build
detailed, non-sentient simulations and so reach the same
trilemma~\cite[\S\,2.3]{Kipping2020bayesian}. Bostrom's FAQ proposes that simulators could
cap simulated civilizations' computing, which would make most civilizations ``leaf nodes
of the simulation tree''~\cite[Q6]{Bostrom2025simulation}.

Bibeau-Delisle and Brassard write a Drake-style equation for the simulated fraction. In
a recursive scenario they bound it above by a simulation-efficiency factor that they
argue ``is most likely under 50\%''~\cite{BibeauDelisle2021probability}.

These authors draw different conclusions from one observation: our civilization has not
produced simulated minds. This comparison is ours, not any single author's. Richmond takes
the observation to make our level unusual on Bostrom's picture, which would favour
hierarchies with fewer levels~\cite{Richmond2008doomsday}. Lewis shows that in some models
conditioning on it \emph{raises} the credence that one is simulated~\cite{Lewis2013doomsday}.
Kipping uses it to obtain a Bayes factor close to unity~\cite{Kipping2020bayesian}.

\subsection{Self-locating belief and anthropic reasoning}
\label{sec:philosophy:anthropics}

The indifference step is one instance of a wider problem: how to assign credence when
one is uncertain which observer one is. Bostrom's \emph{Self-Sampling Assumption} (SSA)
says that one should reason as if one were a random sample from all observers in one's
reference class~\cite[p.~57]{Bostrom2002anthropic}. He later restates it in terms of
observer-moments. The rival \emph{Self-Indication Assumption} (SIA) favours hypotheses on
which many observers exist. Bostrom rejects SIA, mainly through the ``Presumptuous
Philosopher'' case~\cite[pp.~66, 124--125]{Bostrom2002anthropic}. He also shows that
whether an entity belongs to the reference class can change the rational credence when
population size depends on the hypothesis~\cite[pp.~69--71]{Bostrom2002anthropic}.

Elga defends the ``thirder'' answer to Sleeping Beauty~\cite{Elga2000self}. He also
defends an indifference principle: ``similar centered worlds deserve equal
credence''~\cite{Elga2004defeating}. On Elga's view, hypotheses about one's location among
duplicates carry no default bias against them, whereas ordinary sceptical hypotheses
do~\cite{Elga2004defeating}. Birch uses this distinction to specify the rare
circumstances in which Bostrom's selective scepticism could be defensible, and argues that
we are not in them~\cite{Birch2013on}.

Lewis~\cite{Lewis2013doomsday} argues that the simulation argument succeeds where the
Doomsday argument fails. His result depends on a location-uniform prior, which he
identifies with the thirder position. Richmond (abstract-level) argues the opposite:
anti-simulation conclusions follow with ``more \dots\ robust reference classes'' and
without indifference principles~\cite{Richmond2017why}.

Two limitations run through this literature. First, it mostly assumes finitely many
observers~\cite{Bostrom2011patch}. Bostrom's FAQ notes that in an infinite universe the
simulated fraction ``is not defined'', and proposes a limit-density
measure~\cite[Q7]{Bostrom2025simulation}. Second, the verdict depends on the choice of
reference class and self-sampling rule, and that choice is itself disputed.

\subsection{Quantitative treatments}
\label{sec:philosophy:bayes}

Kipping~\cite{Kipping2020bayesian} combines disjuncts (1) and (2) into a ``physical
hypothesis'' $H_P$. He compares it with a simulation hypothesis $H_S$ in which a base
civilization runs $\lambda$ simulations, each of which runs its own simulations with
probability $p$, down to a maximum depth $G$. Assigning equal prior probability to $H_P$
and $H_S$, and conditioning only on one's existence, he obtains a probability of being in
base reality of $\tfrac12 + 1/[2(N_{\mathrm{sim}}+1)]$, where $N_{\mathrm{sim}}$ is the
total number of simulated realities. The probability of being simulated is therefore below
50\% and tends to 50\% as $N_{\mathrm{sim}}\to\infty$. Conditioning instead on the fact
that we have not produced simulations gives a Bayes factor of $(\lambda-1)/\lambda$. Kipping
also finds that if humanity began producing such simulations, the probability of base
reality would fall to $\lambda/N_{\mathrm{sim}}$~\cite{Kipping2020bayesian}.

The 50\% ceiling comes from the equal-prior assumption, not from data. Kipping himself
notes that equal priors could be ``too generous'' to $H_S$ because it is the more complex
model~\cite[\S\,3]{Kipping2020bayesian}. The v1 text also contains two prose sentences that
contradict its own equations; we rely on the abstract and equations. The
Kipping and Bibeau--Brassard models~\cite{Kipping2020bayesian,BibeauDelisle2021probability}
make Bostrom's framework explicit. The numbers they produce follow from assumed priors,
branching structure and efficiency factors, none of which is measured.

\subsection{Simulation and scepticism}
\label{sec:philosophy:scepticism}

Chalmers argues that the hypothesis that one is in a matrix ``is not a skeptical
hypothesis, but a metaphysical hypothesis''~\cite{Chalmers2005matrix}. He treats it as
equivalent to three claims: our world was created, its microphysics is computational, and
our minds are constituted outside it. If this is right, most ordinary beliefs remain true.
Local or recent simulations are at worst ``partial skeptical
hypotheses''~\cite{Chalmers2005matrix}.

In \emph{Reality+}~\cite{Chalmers2022reality}, as summarized in his
précis~\cite{Chalmers2024precis}, Chalmers restates the argument as ``if there are no sim
blockers, we are probably sims''. His sim blockers include Bostrom's two disjuncts and
also the impossibility of intelligent or conscious sims, simulators' avoidance of
conscious sims, and excessive computing cost. He argues that we cannot know that any sim
blocker obtains, so we cannot know that we are not simulated~\cite{Chalmers2024precis}.
His anti-sceptical move is structuralist: a simulation with the same causal structure as
an unsimulated universe makes the same physical theories true~\cite{Chalmers2024precis}.
Replying to critics, he gives an illustration: 10\% credence in the denial of each of two
sim blockers leaves ``a 1\% credence that most humanlike beings are sims''. He rejects
global scepticism but leaves local scepticism open~\cite{Chalmers2024taking}.

Schwitzgebel (abstract-level) argues that, conditional on being simulated, we should give
significant credence to a small or brief simulation~\cite{Schwitzgebel2024lets}.
Summers and Arvan argue that, under panpsychism or panqualityism, viable simulation
hypotheses become ``substantially sceptical''~\cite{Summers2022two}.

\paragraph{Ethics and theology.}
In peer-reviewed philosophy of religion, Dainton proposes a ``simulation solution'' to the
problem of natural evil~\cite{Dainton2020natural}. Crummett assesses it against Fall and
diabolical theodicies~\cite{Crummett2021real}. Johnson argues that some theists are
committed to our world most likely being a simulation~\cite{Johnson2011natural}, and
Steinhart draws out further theological analogues~\cite{Steinhart2010theological}. All four
are cited at abstract level.

\subsection{What the argument does and does not license empirically}
\label{sec:philosophy:bridge}

Five points link this section to the physical proposals reviewed in later sections.

\begin{enumerate}
\item \emph{The argument does not predict any observable signature.} Its conclusion is a
disjunction whose third disjunct concerns the fraction of observers who are simulated,
not the physics they observe~\cite{Bostrom2003are,Bostrom2005simulation}. Bostrom's
resource estimate assumes that simulators render only what observers
notice~\cite[\S\,III]{Bostrom2003are}. His FAQ doubts that lattice-artefact searches of the
kind proposed by Beane \emph{et al.} would detect a simulation. It gives two reasons:
simulators need not use ``the crude simulation technique that such a test would
detect'', and they ``could fake the results''~\cite[Q12]{Bostrom2025simulation}. Any physical test must therefore add
assumptions about how a simulator would be built. Those assumptions are not supplied by
the argument.

\item \emph{Observations bear on it only indirectly.} Bostrom holds that the hypothesis is
testable only in the weak sense that evidence bearing on disjuncts (1) or (2) shifts
credence in~(3)~\cite[Q11]{Bostrom2025simulation}. Qualitative expectations do exist in the
literature. For example, Barrow suggests occasional ``glitches'' and slow drifts in the
constants~\cite{Barrow2007living}. These are motivated guesses, not consequences of the
argument.

\item \emph{Null results and positive results are asymmetric.} Greene argues that no
experiment could show that we are \emph{not} simulated, since a simulator could program
any observation; only a positive anomaly would be diagnostic~\cite{Greene2020termination}.
Bostrom's FAQ makes a related point from the other direction: philosophical arguments, but
not experiments, are immune to the concern that a successful test might prompt
termination~\cite[Q12]{Bostrom2025simulation}.

\item \emph{Physics constrains the premises, not the conclusion.} Empirical bounds on the
cost of simulating physics bear on the empirical premise about computing capacity, that
is, on the ``sims will require too much computer power'' blocker~\cite{Chalmers2024precis}.
But Birch argues that a simulated observer's physics need not reflect the true limits of
computation outside the simulation~\cite{Birch2013on}. Bostrom's FAQ replies with a
two-case argument~\cite[Q4]{Bostrom2025simulation}. Neither side's reasoning is settled by
data.

\item \emph{Credences in the literature rest on priors.} Stated probabilities range from
Bostrom's ``roughly even'' split~\cite{Bostrom2003are}, to ``probably not
simulated''~\cite{Bostrom2009simulation}, to Kipping's ``less than
50\%''~\cite{Kipping2020bayesian}. Each follows from contested indifference principles,
reference classes or model priors. None is a measurement.
\end{enumerate}
